% ===================================================================
%  图表第 12 号 — 车站。— 铁路。— 船。
%  Traduction 2026 d'apres texto/io, controlee sur texto/fr, texto/en
%  et texto/es. Consigne : voir texto/zh/10-tubiao-01.tex.
% ===================================================================

%%K t12-apar-1 apar
\VUsaut{4.0mm}
\VUtitre{15.0pt}{60}{图表第 12 号}
\VUsaut{2.4mm}
\VUpk{11.4pt}{40}{车站。--- 铁路。--- 船。}
\VUsaut{3.4mm}
\textsc{附注。} --- \textit{各国用的时刻不同，所以我们拿}
\VUgras{法国挂图}\textit{上的钟点作例子。}

%%K t12-01-1 p
1. --- 我刚在德国走了一趟。动身以前我买了一本\VUgras{时刻表} \textsuperscript{(15)}，
好知道火车开出的钟点。查明最方便的一班晚上九点离巴黎、一点十三到斯特拉
斯堡以后，我就收拾行装，第二天叫车把我送到车站。

%%K t12-02-1 p
2. --- 我跟在一位乡下老\VUgras{神父} \textsuperscript{(2)} 后面走进大出发厅
时——他提着\VUgras{旅行袋} \textsuperscript{(3)}——已经来了不少\VUgras{旅客}
\textsuperscript{(1)}。

%%K t12-02-2 p
因为我要发一封\VUgras{电报} \textsuperscript{(5)}，就请一位\VUgras{稽查员}
\textsuperscript{(6)} 给我指了\VUgras{电报房} \textsuperscript{(4)}。

%%K t12-03-1 p
3. --- 进\VUgras{候车室} \textsuperscript{(7)} 以前，我先去买了一张来回
\VUgras{车票} \textsuperscript{(9)}。

%%K t12-04-1 p
4. --- 在\VUgras{售票窗} \textsuperscript{(8)} 前我没等多久，因为人不多。一个
休假出门的\VUgras{兵} \textsuperscript{(10)} 好意让我先买，所以我还有工夫向
\VUgras{站长} \textsuperscript{(13)} 问了几句；他正同\VUgras{警察分局长}
\textsuperscript{(12)} 和值班的\VUgras{宪兵} \textsuperscript{(11)} 说话。

%%K t12-tit-1 sub
\VUpk{10.2pt}{40}{托运行李。}
\VUsaut{3.0mm}

%%K t12-05-1 p
5. --- 在\VUgras{书报摊} \textsuperscript{(14)} 买了张报以后，我叫人称我的
\VUgras{行李} \textsuperscript{(17)}——\VUgras{脚夫} \textsuperscript{(16)} 刚用
\VUgras{行李车} \textsuperscript{(19)} 把它推来。\VUgras{办事员} \textsuperscript{(22)} 管着\VUgras{磅秤}
\textsuperscript{(21)}，他说我的箱子重三十五公斤，\VUgras{超重}
五公斤，得在\VUgras{行李窗} \textsuperscript{(23)} 补票钱。

%%K t12-06-1 p
6. --- 我来得早，正好看看站里旅客的来去。有的在\VUgras{寄存处}
\textsuperscript{(25)} 存或取小包；有的在看\VUgras{时刻表} \textsuperscript{(26)}。到站
的旅客往\VUgras{出口} \textsuperscript{(32)} 赶，手里提着\VUgras{手提箱}
\textsuperscript{(24)}。有几个在\VUgras{行李交付处} \textsuperscript{(27)} 等着，递上\VUgras{行李票}
\textsuperscript{(29)} 去领他们的箱子、\VUgras{板条箱} \textsuperscript{(28)} 或
\VUgras{篮子} \textsuperscript{(30)}。

%%K t12-07-1 p
7. --- \VUgras{税吏} \textsuperscript{(33)} 仔细翻检包裹，看有没有该报的东西。

%%K t12-08-1 p
8. --- 有些旅客往\VUgras{餐室} \textsuperscript{(34)} 去，一个\VUgras{侍者}
\textsuperscript{(35)} 忙着招呼他们。他们得快些，因为那趟\VUgras{火车}
\textsuperscript{(36)} 是快车，在站上停不久。

%%K t12-09-1 p
9. --- 随后我走上出发月台，找一节直达\VUgras{车厢} \textsuperscript{(37)}，免得
路上换车。我上了一节有走廊的车，里面有吸烟\VUgras{隔间} \textsuperscript{(38)}，
还有一间留给「单身女客」。

%%K t12-tit-2 sub
\VUpk{10.2pt}{40}{夜里启程。}
\VUsaut{3.0mm}

%%K t12-10-1 p
10. --- 我踏上\VUgras{踏板} \textsuperscript{(40)}，关上\VUgras{车门}
\textsuperscript{(39)}，卷起\VUgras{窗帘} \textsuperscript{(41)}，把小包放上
\VUgras{行李架} \textsuperscript{(43)}，然后在\VUgras{座位} \textsuperscript{(42)} 上
坐下。看见\VUgras{点灯的} \textsuperscript{(44)} 在点灯，我想起要在车里过一夜，
就叫那管出租\VUgras{枕头} \textsuperscript{(45)} 的服务员给我拿一个。

%%K t12-11-1 p
11. --- 查票员来查票。我问他为什么还不开车，时候已经到了。他说
\VUgras{列车长} \textsuperscript{(46)} 正忙着把自行车装进\VUgras{行李车厢}
\textsuperscript{(47)}。再说，\VUgras{暖脚器} \textsuperscript{(48)} 也还没有全换过。

%%K t12-12-1 p
12. --- 不过我们就要开了，因为\VUgras{司炉} \textsuperscript{(50)} 站在他的
\VUgras{煤水车} \textsuperscript{(49)} 上取煤，好添\VUgras{机车} \textsuperscript{(54)}
的火，\VUgras{司机} \textsuperscript{(51)} 只等\VUgras{副站长} \textsuperscript{(52)}
给信号，那副站长站在\VUgras{月台} \textsuperscript{(53)} 上。

%%K t12-13-1 p
13. --- 机车的\VUgras{锅炉} \textsuperscript{(56)} 沉沉地响着；\VUgras{烟囱}
\textsuperscript{(55)} 喷出大股夹着\VUgras{汽} \textsuperscript{(60)} 的烟。一声尖利的
\VUgras{汽笛} \textsuperscript{(59)} 响了，\VUgras{活塞} \textsuperscript{(58)} 动了
起来，带动这沉重机器的\VUgras{轮子} \textsuperscript{(57)}。

%%K t12-tit-3 sub
\VUsaut{3.0mm}
\VUpk{10.2pt}{40}{开车了！}
\VUsaut{3.0mm}

%%K t12-14-1 p
14. --- 火车在\VUgras{铁轨} \textsuperscript{(64)} 上跑起来，越跑越快；
\VUgras{月台} \textsuperscript{(65)} 到了头；我们过了\VUgras{厕所}
\textsuperscript{(66)}，也过了停在另一条\VUgras{线} \textsuperscript{(63)} 上的一列
车厢：\VUgras{卧车} \textsuperscript{(67)}、一节\VUgras{餐车} \textsuperscript{(68)}、
一节\VUgras{邮车} \textsuperscript{(69)}。

%%K t12-noto-1 noto
\VUnotes{143.2mm}{%
\textit{(*) 附注。} --- \textit{这趟旅行自然是按战前的边界情形写的。}%
}

%%K t12-15-1 p
15. --- 接着\VUgras{货站} \textsuperscript{(71)} 从我们眼前过去。棚子底下，职员
们在装慢车托运的\VUgras{货物} \textsuperscript{(72)}。\VUgras{调车工}
\textsuperscript{(76)} 在\VUgras{水塔} \textsuperscript{(73)} 附近调着\VUgras{牲口车}
\textsuperscript{(75)} 和\VUgras{平板车} \textsuperscript{(79)}，编成一列\VUgras{货车}
\textsuperscript{(80)}。

%%K t12-16-1 p
16. --- 我们过了最后一处\VUgras{道岔} \textsuperscript{(78)}，扳道工就站在旁边；
又过了\VUgras{信号机} \textsuperscript{(86)}，车站便消失在远处了。

%%K t12-17-1 p
17. --- 不久我们驶上那座\VUgras{铁桥} \textsuperscript{(74)}，它跨过
\VUgras{河} \textsuperscript{(81)}。过了桥，我们沿河跑；河上有力的\VUgras{拖轮} \textsuperscript{(82)}
拉着沉重的\VUgras{驳船} \textsuperscript{(83)}。我们还看见一条\VUgras{客船}
\textsuperscript{(84)} 刚靠上\VUgras{趸船码头} \textsuperscript{(85)}。

%%K t12-18-1 p
18. --- 飞过几个车站以后，我们钻过几条\VUgras{隧道} \textsuperscript{(87)}，把
数不清的\VUgras{道口看守}的\VUgras{小屋} \textsuperscript{(88)} 抛在后面。给车身
的摇晃摇着，我沉沉睡去。

%%K t12-tit-4 sub
\VUsaut{3.0mm}
\VUpk{10.2pt}{40}{多伊奇－阿夫里库尔站。}
\VUsaut{3.0mm}

%%K t12-19-1 p
19. --- 我忽然被一个职员的喊声惊醒：「多伊奇－阿夫里库尔！
\textsuperscript{(*)} 全体换车！」那是海关检查和边界上一切啰唆的手续。我得照
站上的\VUgras{钟} \textsuperscript{(89)} 对表，那钟快了五十五分钟；我睡眼惺忪，
几乎分不清\VUgras{长针} \textsuperscript{(90)} 和\VUgras{短针} \textsuperscript{(91)}；
连\VUgras{钟面} \textsuperscript{(92)} 也叫早晨的雾模糊了。

%%K t12-20-1 p
20. --- 海关的人检查的工夫，我打着呵欠看那些贴在站墙上的\VUgras{招贴}
\textsuperscript{(93)}。我吃了早饭消磨时间，又照德国的\VUgras{路网} \textsuperscript{(94)}
图看了看离斯特拉斯堡还有多远，然后回到自己的车厢，心里因为快到目的地
而振作起来。
\VUsaut{6.0mm}
\VUfilet{22mm}
